% `tex.web` §392 and §396: a `\par` in a macro's argument, and the one place it
% is not a fault.
%
%     continue: get_token;
%     if cur_tok=info(r) then
%       @<Advance r; goto found if the parameter delimiter has been fully
%         matched, otherwise goto continue@>;
%     @<Contribute the recently matched tokens to the current parameter …@>;
%     if cur_tok=par_token then if long_state<>long_call then
%       @<Report a runaway argument and abort@>;
%
% The delimiter is tested first. So `\def\a#1\par#2{}` — a macro whose first
% argument ends at a paragraph — takes a `\par` there without complaint, while
% the same token anywhere else in a non-`\long` argument is "Paragraph ended
% before \a was complete".
%
% We tested for `\par` first and abandoned every call of such a macro. trip.tex
% line 354 defines one on purpose, and the run stopped there — with the
% argument scan running away to the end of the file, which is the very thing
% §396 exists to prevent.
\catcode`\{=1 \catcode`\}=2 \catcode`\#=6

% `\par` as the delimiter: no error, and it is not part of the argument.
\def\a#1\par#2{[#1|#2]}
\message{PROBE:one=\a x\par y}

% `\par` inside the argument of a delimiter that is not `\par`: §396's error,
% and the call is abandoned.
\def\b#1x{[#1]}
\message{PROBE:two=before}
\b a\par x
\message{PROBE:three=after}

% ...and `\long` allows it.
\long\def\c#1x{[#1]}
\message{PROBE:four=\c a\par x}
\message{PROBE:done}
\end
